\chapter{Multiple and Logistic Regression}
\label{ch:multiple-logistic-regression}

Chapter~12 fitted a straight line through a cloud of points: one predictor $x$, one
response $y$. Real questions rarely have one cause. A house price depends on size
\emph{and} age \emph{and} location. A page conversion depends on time on page
\emph{and} device \emph{and} traffic source. This chapter does two things. It lets the
line take in many predictors at once, which turns it into a plane and then into a
projection in high-dimensional space. Then it bends the model so it can predict a
yes/no outcome. The second step leads straight to the loss function that trains most
classifiers in machine learning.

\section{From a line to a plane}
\label{sec:mlr-plane}

Here are the first rows of a dataset of $40$ houses. Price is in thousands of dollars,
size in square metres, age in years.

\begin{center}
\begin{tabular}{rrrr}
\toprule
house & size (m$^2$) & age (yr) & price (\$k) \\
\midrule
1 & 144 & 16 & 308 \\
2 & 185 & 53 & 343 \\
3 & 166 & 31 & 313 \\
4 & 84 & 51 & 161 \\
5 & 95 & 38 & 142 \\
$\vdots$ & $\vdots$ & $\vdots$ & $\vdots$ \\
\bottomrule
\end{tabular}
\end{center}

With one predictor we drew $\hat y = b_0 + b_1 x$, a line. With two predictors the
natural guess is
\[
  \widehat{\text{price}} = b_0 + b_1\cdot\text{size} + b_2\cdot\text{age},
\]
which is a flat plane hovering over the (size, age) floor. Least squares again picks
the plane that makes the sum of squared vertical gaps as small as possible. For these
$40$ houses the computer returns
\[
  \widehat{\text{price}} = 60.08 + 1.861\cdot\text{size} - 1.393\cdot\text{age}.
\]
Read it as: one extra square metre adds about \$1{,}861; one extra year of age
subtracts about \$1{,}393.

\begin{figure}[H]
\centering
\begin{tikzpicture}
\begin{axis}[
  width=11cm, height=8cm, view={-35}{22},
  xlabel={size (m$^2$)}, ylabel={age (yr)}, zlabel={price (\$k)},
  xmin=50, xmax=200, ymin=0, ymax=60, zmin=50, zmax=420,
  xtick={100,150,200}, ytick={20,40,60},
  label style={font=\small}, tick label style={font=\footnotesize},
  grid=major,
  legend style={font=\footnotesize, at={(0.5,1.03)}, anchor=south, legend columns=3,
    /tikz/every even column/.append style={column sep=0.3cm}},
]
\addplot3[surf, domain=50:200, y domain=0:60, samples=6, samples y=5,
  opacity=0.35, fill opacity=0.35, colormap={blues}{color=(histfill) color=(sourceblue)},
  faceted color=sourceblue!60, forget plot] {60.08+1.8615*x-1.3934*y};
\draw[residgray, very thick] (axis cs:144,16,308) -- (axis cs:144,16,305.8);
\draw[residgray, very thick] (axis cs:185,53,343) -- (axis cs:185,53,330.6);
\draw[residgray, very thick] (axis cs:166,31,313) -- (axis cs:166,31,325.9);
\draw[residgray, very thick] (axis cs:84,51,161) -- (axis cs:84,51,145.4);
\draw[residgray, very thick] (axis cs:95,38,142) -- (axis cs:95,38,184.0);
\draw[residgray, very thick] (axis cs:181,45,314) -- (axis cs:181,45,334.3);
\draw[residgray, very thick] (axis cs:51,5,153) -- (axis cs:51,5,148.0);
\draw[residgray, very thick] (axis cs:173,32,287) -- (axis cs:173,32,337.5);
\draw[residgray, very thick] (axis cs:170,30,358) -- (axis cs:170,30,334.7);
\draw[residgray, very thick] (axis cs:120,52,184) -- (axis cs:120,52,211.0);
\draw[residgray, very thick] (axis cs:95,22,230) -- (axis cs:95,22,206.3);
\draw[residgray, very thick] (axis cs:92,36,168) -- (axis cs:92,36,181.2);
\draw[residgray, very thick] (axis cs:88,4,232) -- (axis cs:88,4,218.3);
\draw[residgray, very thick] (axis cs:117,23,250) -- (axis cs:117,23,245.8);
\addlegendimage{residgray, very thick}
\addlegendentry{residual (vertical gap to plane)}
\addplot3[only marks, mark=*, mark size=2pt, densityred] coordinates {
(144,16,308) (185,53,343) (84,51,161) (51,5,153) (170,30,358) (95,22,230) (88,4,232) (117,23,250)};
\addlegendentry{house above the plane}
\addplot3[only marks, mark=square*, mark size=1.8pt, sourceblue!80!black] coordinates {
(166,31,313) (95,38,142) (181,45,314) (173,32,287) (120,52,184) (92,36,168)};
\addlegendentry{house below the plane}
\end{axis}
\end{tikzpicture}
\caption{With two predictors the fitted model is a plane: the predicted price for every
(size, age) pair. Red dots are houses priced above the plane, blue squares houses priced
below it, and each grey stick is a residual. Least squares tilts the plane to make the sum
of the squared stick lengths as small as possible, exactly as in Chapter~12. (Fourteen of
the $40$ houses are shown.)}
\label{fig:mlr-plane}
\end{figure}

With three or more predictors there is no picture of the plane any more, but the
algebra does not care. We need a notation that handles any number of predictors at once.

\section{The matrix form}
\label{sec:mlr-matrix}

\subsection{Stacking the equations}

Take a tiny dataset so we can do everything by hand: four houses, size $x$ in units of
$100$ m$^2$, price $y$ in units of \$100k.
\[
  x = (1,\,2,\,3,\,4), \qquad y = (3,\,5,\,4,\,8).
\]
The model says each price is intercept plus slope times size plus an error:
\[
\begin{aligned}
  3 &= \beta_0 + \beta_1\cdot 1 + \eps_1,\\
  5 &= \beta_0 + \beta_1\cdot 2 + \eps_2,\\
  4 &= \beta_0 + \beta_1\cdot 3 + \eps_3,\\
  8 &= \beta_0 + \beta_1\cdot 4 + \eps_4.
\end{aligned}
\]
Four equations with the same shape. Stack them into vectors and a matrix:
\[
  \underbrace{\begin{bmatrix}3\\5\\4\\8\end{bmatrix}}_{y}
  =
  \underbrace{\begin{bmatrix}1&1\\1&2\\1&3\\1&4\end{bmatrix}}_{X}
  \underbrace{\begin{bmatrix}\beta_0\\ \beta_1\end{bmatrix}}_{\beta}
  +
  \underbrace{\begin{bmatrix}\eps_1\\ \eps_2\\ \eps_3\\ \eps_4\end{bmatrix}}_{\eps}.
\]

\begin{definition}[Linear model in matrix form]
With $n$ observations and $p$ predictors,
\[
  y = X\beta + \eps,
\]
where $y$ is the $n\times 1$ vector of responses, $\beta$ is the $(p+1)\times 1$ vector of
coefficients, $\eps$ is the $n\times 1$ vector of errors, and $X$ is the $n\times(p+1)$
\textbf{design matrix}: one row per observation, one column per predictor, plus a
leading column of $1$s.
\end{definition}

The design matrix is the spreadsheet of inputs. The column of $1$s is the
\textbf{intercept column}: it multiplies $\beta_0$, so every row gets the same constant
added. Without it the fitted plane would be forced through the origin.

\begin{figure}[H]
\centering
\begin{tikzpicture}[font=\small]
\node (y) at (0,0) {$\begin{bmatrix} y_1\\ y_2\\ \vdots\\ y_n \end{bmatrix}$};
\node at (0.95,0) {$=$};
\matrix (M) at (4.1,0) [matrix of math nodes, left delimiter={[}, right delimiter={]},
  nodes={minimum width=0.75cm, minimum height=0.55cm, anchor=center}, column sep=1pt, row sep=0pt]
{ 1 & x_{11} & x_{12} & \cdots & x_{1p}\\
  1 & x_{21} & x_{22} & \cdots & x_{2p}\\
  \vdots & \vdots & \vdots & & \vdots\\
  1 & x_{n1} & x_{n2} & \cdots & x_{np}\\ };
\node (b) at (7.55,0) {$\begin{bmatrix} \beta_0\\ \beta_1\\ \vdots\\ \beta_p \end{bmatrix}$};
\node at (8.35,0) {$+$};
\node (e) at (9.15,0) {$\begin{bmatrix} \eps_1\\ \eps_2\\ \vdots\\ \eps_n \end{bmatrix}$};
\node[draw=quotientorange, thick, rounded corners, inner sep=1pt, fit=(M-1-1)(M-4-1)] (IC) {};
\node[draw=sourceblue, thick, rounded corners, inner sep=1pt, fit=(M-1-2)(M-4-5)] (PC) {};
\node[draw=targetgreen, thick, rounded corners, inner sep=0pt, fit=(M-2-1)(M-2-5)] (R2) {};
\node[quotientorange, anchor=north, align=center] at (IC.south) {intercept\\column};
\node[sourceblue, anchor=north, align=center] at (PC.south) {one column\\per predictor};
\node[targetgreen, anchor=south] at (PC.north) {one row per observation (row 2 shown)};
\node[anchor=north, align=center] at (0,-1.15) {$n\times 1$\\responses};
\node[anchor=north, align=center] at (7.55,-1.15) {$(p{+}1)\times 1$\\coefficients};
\node[anchor=north, align=center] at (9.3,-1.15) {$n\times 1$\\errors};
\end{tikzpicture}
\caption{The anatomy of $y = X\beta + \eps$. The design matrix $X$ holds the inputs;
its first column of $1$s carries the intercept. Multiplying row $i$ of $X$ by $\beta$
gives the model's prediction for observation $i$.}
\label{fig:mlr-design}
\end{figure}

\subsection{The normal equations}

We want the $\beta$ that minimises the sum of squared residuals
\[
  S(\beta) = \sum_{i=1}^n (y_i - \hat y_i)^2 = \norm{y - X\beta}^2 .
\]
The symbol $\norm{v}^2$ is the squared length of a vector: the sum of its squared
entries. So $S(\beta)$ is the squared distance between the data vector $y$ and the
prediction vector $X\beta$.

Expand it, treating vectors like numbers but keeping track of transposes
($v^\top$ turns a column into a row):
\[
  S(\beta) = y^\top y - 2\beta^\top X^\top y + \beta^\top X^\top X \beta .
\]
This is a bowl-shaped quadratic in $\beta$. At the bottom of a bowl the slope in every
direction is zero. The vector of slopes (the \textbf{gradient}, one partial
derivative per coefficient) is
\[
  \nabla S(\beta) = -2X^\top y + 2X^\top X\beta .
\]
Setting it to zero gives the central equation of this chapter.

\begin{theorem}[Normal equations]
The least squares coefficients $\hat\beta = (b_0, b_1, \dots, b_p)^\top$ satisfy
\[
  X^\top X\,\hat\beta = X^\top y .
\]
If the columns of $X$ are linearly independent, $X^\top X$ is invertible and
\[
  \hat\beta = (X^\top X)^{-1} X^\top y .
\]
\end{theorem}

\begin{example}[Four houses by hand]
\label{ex:mlr-four}
With $X$ and $y$ above,
\[
  X^\top X = \begin{bmatrix} 4 & 10\\ 10 & 30 \end{bmatrix},\qquad
  X^\top y = \begin{bmatrix} 3+5+4+8\\ 3+10+12+32 \end{bmatrix}
           = \begin{bmatrix} 20\\ 57 \end{bmatrix}.
\]
The entries of $X^\top X$ are $n=4$, $\sum x_i = 10$ and $\sum x_i^2 = 30$. The
determinant is $4\cdot 30 - 10\cdot 10 = 20$, so
\[
  \hat\beta = \frac{1}{20}\begin{bmatrix} 30 & -10\\ -10 & 4 \end{bmatrix}
  \begin{bmatrix} 20\\ 57 \end{bmatrix}
  = \frac{1}{20}\begin{bmatrix} 600-570\\ -200+228 \end{bmatrix}
  = \begin{bmatrix} 1.5\\ 1.4 \end{bmatrix}.
\]
The fitted line is $\hat y = 1.5 + 1.4x$, the same answer the slope formula
$b_1 = S_{xy}/S_{xx}$ from Chapter~12 gives. Fitted values are $2.9, 4.3, 5.7, 7.1$ and
residuals are $e = (0.1,\ 0.7,\ -1.7,\ 0.9)$.
\end{example}

Two checks on those residuals. They sum to zero: $0.1+0.7-1.7+0.9=0$. And weighted by
$x$ they also sum to zero: $0.1 + 1.4 - 5.1 + 3.6 = 0$. Those two facts are the normal
equations rewritten as $X^\top e = 0$: the residual vector has zero dot product with
every column of $X$. That sentence has a picture.

\begin{lstlisting}[language=Python]
import numpy as np
X = np.array([[1, 1], [1, 2], [1, 3], [1, 4]], dtype=float)
y = np.array([3, 5, 4, 8], dtype=float)
beta = np.linalg.solve(X.T @ X, X.T @ y)       # normal equations: [1.5, 1.4]
beta2, *_ = np.linalg.lstsq(X, y, rcond=None)  # same answer, numerically safer
e = y - X @ beta
print(beta, X.T @ e)                           # X^T e is (0, 0) up to rounding
\end{lstlisting}

\section{The geometric picture: least squares is a projection}
\label{sec:mlr-geometry}

Think of the $n$ responses as a single arrow $y$ in $n$-dimensional space. Each column of
$X$ is also an arrow in that space. Every possible prediction $X\beta$ is a weighted
sum of the columns, so the set of all possible predictions is a flat sheet through the
origin: the \textbf{column space} of $X$, meaning every vector you can build by
stretching and adding the columns. With two columns it is a plane.

The data arrow $y$ almost never lies in that plane (noise pushes it out). Least squares
asks: which point in the plane is closest to $y$? The answer is the foot of the
perpendicular dropped from $y$ onto the plane. That foot is $\hat y$, and the
perpendicular itself is the residual vector $e = y - \hat y$.

\begin{figure}[H]
\centering
\begin{tikzpicture}[scale=1.15, font=\small]
\fill[histfill!25] (-0.8,-1.0) -- (5.4,-1.0) -- (6.8,1.6) -- (0.6,1.6) -- cycle;
\draw[histfill!80!black] (-0.8,-1.0) -- (5.4,-1.0) -- (6.8,1.6) -- (0.6,1.6) -- cycle;
\node[sourceblue!80!black, anchor=north west] at (-0.8,-1.05) {column space of $X$};
\coordinate (O) at (0.6,-0.4);
\coordinate (Yh) at (3.9,0.45);
\coordinate (Y) at (3.9,3.3);
\coordinate (Z) at (2.9,1.2);
\draw[-{Stealth}, thick, gray!70!black] (O) -- (2.4,0.55) node[right] {$\mathbf{1}$};
\draw[-{Stealth}, thick, gray!70!black] (O) -- (4.6,-0.55) node[below] {$x$};
\draw[-{Stealth}, very thick, distblue] (O) -- (Y) node[above] {$y$ (data)};
\draw[-{Stealth}, very thick, targetgreen] (O) -- (Yh);
\node[targetgreen!70!black, anchor=north west] at (Yh) {$\hat y = X\hat\beta$};
\draw[-{Stealth}, very thick, densityred, dashed] (Yh) -- (Y);
\node[densityred, anchor=west] at (3.95,1.9) {$e = y-\hat y$};
\draw[densityred] ($(Yh)+(0,0.28)$) -- ($(Yh)+(0.26,0.36)$) -- ($(Yh)+(0.26,0.08)$);
\draw[trapred, dotted, thick] (Y) -- (Z);
\fill[trapred] (Z) circle (1.4pt);
\node[trapred, anchor=north, font=\footnotesize] at (2.9,1.13) {another fit $Xb$};
\end{tikzpicture}
\caption{Least squares as projection. The grey arrows are the two columns of $X$: the
column of $1$s ($\mathbf 1$) and the predictor column ($x$). The data vector $y$ sticks out
of the plane of possible predictions they span. The closest point in the plane is $\hat y$, reached by dropping a
perpendicular (the small square marks the right angle). Any other fit $Xb$ in the plane
(dotted) is farther from $y$. The residual $e$ is at a right angle to every column of $X$,
which is exactly the statement $X^\top e = 0$.}
\label{fig:mlr-projection}
\end{figure}

\begin{whybox}
Why does the perpendicular give the closest point? Take any other point $Xb$ in the
plane. The three points $y$, $\hat y$, $Xb$ form a right triangle with the right angle
at $\hat y$, because $\hat y - Xb$ lies in the plane and $e$ is perpendicular to the
plane. Pythagoras gives
\[
  \norm{y - Xb}^2 = \norm{e}^2 + \norm{\hat y - Xb}^2 \;\ge\; \norm{e}^2 .
\]
So no other point beats $\hat y$. Writing ``$e$ is perpendicular to every column'' as
$X^\top(y - X\hat\beta) = 0$ and rearranging gives the normal equations again, with no
calculus at all.
\end{whybox}

The same right triangle, with $\bar y\mathbf 1$ in place of $Xb$, is the decomposition
$\text{SST} = \text{SSR} + \text{SSE}$ from Chapter~12. $R^2$ is the squared cosine of the
angle between the centred $y$ and the plane: how much of the arrow lies \emph{in} the
plane.

\subsection{The hat matrix}

Substitute $\hat\beta$ into $\hat y = X\hat\beta$:
\[
  \hat y = X(X^\top X)^{-1}X^\top y = Hy, \qquad H = X(X^\top X)^{-1}X^\top .
\]
$H$ is called the \textbf{hat matrix} because it puts the hat on $y$. It is the
projection machine: feed in any vector, get back its shadow on the column space.
Two facts follow from the picture. Projecting twice changes nothing, so $HH = H$. And the
residual is what is left over, $e = (I - H)y$.

The diagonal entry $h_{ii}$ is the \textbf{leverage} of observation $i$: how strongly
$y_i$ pulls on its own fitted value $\hat y_i$. For the four houses, with one predictor,
$h_{ii} = \tfrac1n + (x_i-\xbar)^2/S_{xx}$ with $\xbar = 2.5$ and $S_{xx} = 5$:
\[
  h_{11} = 0.25 + \tfrac{2.25}{5} = 0.7,\quad h_{22} = 0.25 + \tfrac{0.25}{5} = 0.3,\quad
  h_{33} = 0.3,\quad h_{44} = 0.7 .
\]
The end points have the most leverage, as the seesaw intuition of Chapter~12 suggested.
The leverages sum to $2$, the number of columns. In general they sum to $p+1$.

\begin{recoverybox}
Pause and think: if you add a predictor that is pure noise, can $\norm{e}^2$ go up?
(No. The column space only grows, so the closest point can only get closer. That is why
$R^2$ never falls when a variable is added, and why we will need adjusted $R^2$.)
\end{recoverybox}

\section{What a coefficient means: ``holding the others fixed''}
\label{sec:mlr-interpret}

In $\widehat{\text{price}} = 60.08 + 1.861\cdot\text{size} - 1.393\cdot\text{age}$,
compare two houses with the same age whose sizes differ by $1$ m$^2$. Their predicted
prices differ by exactly $1.861$ (thousand dollars). That is the meaning of a
coefficient in multiple regression: \emph{the predicted change in $y$ for a one-unit
change in $x_j$ with every other predictor held fixed.}

This is a different number from the simple-regression slope, and the difference can be
dramatic. Here are $12$ houses with size and number of bedrooms.

\begin{center}
\begin{tabular}{l*{12}{r}}
\toprule
size (m$^2$) & 60 & 70 & 80 & 90 & 100 & 110 & 120 & 130 & 140 & 150 & 160 & 170\\
bedrooms & 1 & 2 & 1 & 2 & 3 & 2 & 3 & 4 & 3 & 4 & 5 & 4\\
price (\$k) & 161 & 130 & 188 & 185 & 191 & 228 & 219 & 238 & 268 & 307 & 287 & 317\\
\bottomrule
\end{tabular}
\end{center}

Regress price on bedrooms alone:
\[
  \widehat{\text{price}} = 122.6 + 36.7\cdot\text{bedrooms}.
\]
Each extra bedroom is worth \$36{,}700. Now add size:
\[
  \widehat{\text{price}} = 31.5 + 2.127\cdot\text{size} - 17.5\cdot\text{bedrooms}.
\]
The bedroom coefficient flipped sign. Holding size fixed, an extra bedroom
\emph{lowers} the price by \$17{,}500.

\begin{figure}[H]
\centering
\begin{tikzpicture}
\begin{axis}[
  width=11cm, height=7cm,
  xlabel={bedrooms}, ylabel={price (\$k)},
  xmin=0.5, xmax=6.1, ymin=100, ymax=350, xtick={1,2,3,4,5},
  legend style={font=\footnotesize, at={(0.5,-0.2)}, anchor=north, legend columns=2},
  axis lines=left,
]
\addplot[only marks, mark=*, mark size=2.2pt, histfill!70!black] coordinates
  {(1,161) (2,130) (1,188) (2,185) (3,191)};
\addlegendentry{size 60--100 m$^2$}
\addplot[only marks, mark=square*, mark size=2.2pt, quotientorange] coordinates
  {(2,228) (3,219) (4,238)};
\addlegendentry{size 110--130 m$^2$}
\addplot[only marks, mark=triangle*, mark size=2.8pt, densityred] coordinates
  {(3,268) (4,307) (5,287) (4,317)};
\addlegendentry{size 140--170 m$^2$}
\addplot[thick, gray!70!black, domain=0.8:5.2] {122.63+36.689*x};
\addlegendentry{bedrooms only: slope $+36.7$}
\addplot[thick, dashed, histfill!70!black, domain=1:3] {201.68-17.483*x};
\addlegendentry{size fixed: slope $-17.5$}
\addplot[thick, dashed, quotientorange, domain=2:4, forget plot] {286.75-17.483*x};
\addplot[thick, dashed, densityred, domain=3:5, forget plot] {361.18-17.483*x};
\node[font=\footnotesize, histfill!70!black, anchor=west] at (axis cs:3.05,149.2) {size fixed at $80$ m$^2$};
\node[font=\footnotesize, quotientorange, anchor=west] at (axis cs:4.05,216.8) {at $120$ m$^2$};
\node[font=\footnotesize, densityred, anchor=west] at (axis cs:5.05,273.8) {at $155$ m$^2$};
\end{axis}
\end{tikzpicture}
\caption{Confounding in a regression. Across all houses, more bedrooms go with higher
prices (solid grey line, slope $+36.7$), because big houses have more bedrooms. Now hold
size fixed: each dashed line is the multiple-regression prediction at one size ($80$,
$120$ or $155$ m$^2$), and each slopes \emph{down} ($-17.5$ per bedroom), because at a
given size more bedrooms means smaller rooms. Marker shape shows each house's size group.}
\label{fig:mlr-confound}
\end{figure}

\begin{figure}[H]
\centering
\begin{tikzpicture}[font=\small, node distance=2.6cm,
  box/.style={draw, rounded corners, thick, minimum width=2.3cm, minimum height=0.8cm, align=center}]
\node[box, draw=structurepurple] (S) {size};
\node[box, draw=sourceblue, below left=1.2cm and 2.2cm of S] (B) {bedrooms};
\node[box, draw=densityred, below right=1.2cm and 2.2cm of S] (P) {price};
\draw[-{Stealth}, thick] (S) -- node[above left, align=center] {$\delta$: houses with one more\\bedroom are $25.5$ m$^2$\\bigger on average} (B);
\draw[-{Stealth}, thick] (S) -- node[above right, align=center] {$+\$2.13$k per m$^2$\\(bedrooms fixed)} (P);
\draw[-{Stealth}, thick] (B) -- node[below=3pt, align=center] {direct effect: $-\$17.5$k per bedroom\\(size fixed)} (P);
\end{tikzpicture}
\caption{Why the coefficient flips. Size drives both bedroom count and price. The
bedrooms-only slope adds two routes from bedrooms to price: the direct arrow
($-17.5$) plus the detour through size ($25.5 \times 2.13 \approx 54.2$), giving
$-17.5 + 54.2 = +36.7$.}
\label{fig:mlr-dag}
\end{figure}

The arithmetic in Figure~\ref{fig:mlr-dag} is exact, not a heuristic. Regress size on
bedrooms and you get slope $\delta = 25.47$: each extra bedroom comes with $25.47$ m$^2$
more house on average. Then
\[
  \underbrace{36.69}_{\text{simple slope}}
  = \underbrace{-17.48}_{\text{bedrooms, size fixed}}
  + \underbrace{2.1267}_{\text{size coef}}\times\underbrace{25.47}_{\delta}.
\]
This identity is called the \textbf{omitted variable bias} formula: leaving out a
variable that is correlated with an included one shifts the included coefficient by
(the omitted variable's coefficient) $\times$ ($\delta$, the slope from regressing the
omitted variable on the included one).
It is the regression version of the confounding story from Chapter~5.

\begin{trapbox}
``Holding others fixed'' is a statement about the fitted model, not about the world.
Regression adjusts only for variables you put in. If an unmeasured variable (school
district, say) drives both bedrooms and price, the bedroom coefficient still carries its
shadow. And if no house in your data has $5$ bedrooms in $60$ m$^2$, the model's
prediction for one is extrapolation, not evidence.
\end{trapbox}

\section{Building richer models with the same machinery}
\label{sec:mlr-features}

The word ``linear'' in linear regression means linear in the coefficients $\beta$. The
columns of $X$ can be anything you compute from the data. That freedom gives
categorical predictors, interactions and curves at no extra cost.

\subsection{Categorical predictors and dummy variables}

Request latency (ms) from three regions, four requests each:

\begin{center}
\begin{tabular}{lrrrrr}
\toprule
region & \multicolumn{4}{c}{latencies} & mean \\
\midrule
US & 42 & 45 & 40 & 47 & 43.5\\
EU & 51 & 55 & 49 & 53 & 52.0\\
Asia & 78 & 72 & 75 & 79 & 76.0\\
\bottomrule
\end{tabular}
\end{center}

Region is a label, not a number. Coding US$=1$, EU$=2$, Asia$=3$ would force the EU to sit
exactly halfway between the others. Instead create \textbf{dummy variables}: columns of
$0$s and $1$s that switch on for one category. Pick US as the \textbf{baseline}
(reference category) and add two columns:
\[
  d_{\text{EU}} = \begin{cases}1 & \text{EU}\\0&\text{otherwise}\end{cases},\qquad
  d_{\text{Asia}} = \begin{cases}1 & \text{Asia}\\0&\text{otherwise}\end{cases}.
\]
The model $\hat y = b_0 + b_1 d_{\text{EU}} + b_2 d_{\text{Asia}}$ predicts $b_0$ for
US, $b_0+b_1$ for EU and $b_0 + b_2$ for Asia. Least squares makes each prediction the
group mean, so
\[
  b_0 = 43.5,\qquad b_1 = 52.0 - 43.5 = 8.5,\qquad b_2 = 76.0 - 43.5 = 32.5 .
\]
Each dummy coefficient is a difference from the baseline. Its $t$-test asks whether that
region differs from US. That is the pooled two-sample $t$-test of Chapter~9 in regression
clothing, except that the spread is pooled across all three regions.

\begin{figure}[H]
\centering
\begin{tikzpicture}
\begin{axis}[
  width=10cm, height=6.5cm,
  ylabel={latency (ms)}, xmin=0.4, xmax=4.2, ymin=35, ymax=85,
  xtick={1,2,3}, xticklabels={US (baseline), EU, Asia},
  axis lines=left,
]
\addplot[only marks, mark=*, mark size=2pt, sourceblue] coordinates
  {(0.9,42) (1.0,45) (1.1,40) (0.95,47) (1.9,51) (2.0,55) (2.1,49) (1.95,53)
   (2.9,78) (3.0,72) (3.1,75) (2.95,79)};
\draw[thick, gray!70!black, dashed] (axis cs:0.6,43.5) -- (axis cs:3.6,43.5);
\node[anchor=north west, font=\footnotesize, gray!70!black] at (axis cs:1.3,43.2) {$b_0=43.5$ (US mean)};
\draw[very thick, targetgreen] (axis cs:1.75,52) -- (axis cs:2.25,52);
\draw[very thick, targetgreen] (axis cs:2.75,76) -- (axis cs:3.25,76);
\draw[very thick, targetgreen] (axis cs:0.75,43.5) -- (axis cs:1.25,43.5);
\node[anchor=east, font=\footnotesize, targetgreen!70!black] at (axis cs:1.72,52) {EU mean $52.0$};
\node[anchor=east, font=\footnotesize, targetgreen!70!black] at (axis cs:2.72,76) {Asia mean $76.0$};
\draw[-{Stealth}, thick, densityred] (axis cs:2.35,43.5) -- (axis cs:2.35,52);
\node[anchor=west, font=\footnotesize, densityred] at (axis cs:2.37,47.5) {$b_1=8.5$};
\draw[-{Stealth}, thick, densityred] (axis cs:3.35,43.5) -- (axis cs:3.35,76);
\node[anchor=west, font=\footnotesize, densityred] at (axis cs:3.37,60) {$b_2=32.5$};
\end{axis}
\end{tikzpicture}
\caption{Dummy coding. Blue dots are the $12$ requests; green bars are the three group
means. The intercept $b_0$ is the baseline (US) mean, extended as the dashed line, and each
dummy coefficient is the red jump from that line up to its own group's mean.}
\label{fig:mlr-dummy}
\end{figure}

\begin{trapbox}
The \textbf{dummy variable trap}: with an intercept, use $k-1$ dummies for $k$
categories. If you add $d_{\text{US}}$ too, then $d_{\text{US}} + d_{\text{EU}} +
d_{\text{Asia}} = 1$ in every row, which equals the intercept column. The columns are
linearly dependent, $X^\top X$ has no inverse, and the coefficients are not defined.
Software either errors or silently drops a column.
\end{trapbox}

\subsection{Interactions}

Latency also depends on payload size. Suppose an old and a new server version both get
slower with payload. Two very different stories are possible, and Figure~\ref{fig:mlr-interaction}
shows both. Let $x$ be payload (KB) and $d=1$ for the new version.
\begin{itemize}
  \item \emph{Parallel lines}: $\hat y = 20 + 0.5x - 8d$. The new version is $8$ ms
  faster at every payload.
  \item \emph{Non-parallel lines}: $\hat y = 20 + 0.5x - 8d - 0.2\,xd$. The new version
  is $8$ ms faster at $x=0$ and its slope is $0.5 - 0.2 = 0.3$ ms per KB, so the gap grows
  with payload.
\end{itemize}
The extra column $x\cdot d$ is an \textbf{interaction term}: a product of two predictors,
which lets the effect of one depend on the value of the other.

\begin{figure}[H]
\centering
\begin{tikzpicture}
\begin{axis}[
  name=left, width=6.4cm, height=5.6cm,
  title={no interaction}, title style={font=\small},
  xlabel={payload (KB)}, ylabel={latency (ms)},
  xmin=0, xmax=100, ymin=0, ymax=75, axis lines=left,
]
\addplot[thick, densityred, domain=0:100] {20+0.5*x};
\addplot[thick, sourceblue, domain=0:100] {12+0.5*x};
\node[densityred, font=\footnotesize, anchor=south east] at (axis cs:95,68) {old};
\node[sourceblue, font=\footnotesize, anchor=north west] at (axis cs:62,40) {new};
\draw[{Stealth}-{Stealth}, gray!70!black] (axis cs:20,22.5) -- (axis cs:20,29.5);
\node[gray!70!black, font=\footnotesize, anchor=south] at (axis cs:20,36) {gap $8$};
\end{axis}
\begin{axis}[
  at={(left.east)}, anchor=west, xshift=1.4cm,
  width=6.4cm, height=5.6cm,
  title={with interaction}, title style={font=\small},
  xlabel={payload (KB)},
  xmin=0, xmax=100, ymin=0, ymax=75, axis lines=left,
]
\addplot[thick, densityred, domain=0:100] {20+0.5*x};
\addplot[thick, sourceblue, domain=0:100] {12+0.3*x};
\node[densityred, font=\footnotesize, anchor=south east] at (axis cs:95,68) {old};
\node[sourceblue, font=\footnotesize, anchor=north west] at (axis cs:62,30) {new};
\draw[{Stealth}-{Stealth}, gray!70!black] (axis cs:20,18.5) -- (axis cs:20,29.5);
\node[gray!70!black, font=\footnotesize, anchor=south] at (axis cs:20,36) {gap $12$};
\draw[{Stealth}-{Stealth}, gray!70!black] (axis cs:90,39.5) -- (axis cs:90,64.5);
\node[gray!70!black, font=\footnotesize, anchor=east] at (axis cs:89,50) {gap $26$};
\end{axis}
\end{tikzpicture}
\caption{Predicted latency for the old (red) and new (blue) server versions. Left: without an
interaction term the model can only shift the line, so the gap is the same $8$ ms at every
payload. Right: with an $x\cdot d$ term the slopes may differ, and the gap grows from
$8$ ms at $0$ KB ($12$ ms at $20$ KB, $26$ ms at $90$ KB) to $28$ ms at $100$ KB.}
\label{fig:mlr-interaction}
\end{figure}

\begin{trapbox}
Once an interaction is in the model, the coefficient on $d$ alone ($-8$) is the effect of
the new version \emph{at $x=0$} only. Reporting it as ``the effect of the new version''
is wrong whenever $x=0$ is far from typical payloads. Centring $x$ (subtracting its mean)
before forming the product makes that coefficient the effect at an average payload.
\end{trapbox}

\subsection{Polynomial features and overfitting}

A curve is a line in disguise. To fit $\hat y = b_0 + b_1x + b_2x^2 + b_3x^3$, put the
columns $1, x, x^2, x^3$ into $X$ and run the same normal equations. The model is
nonlinear in $x$ but linear in the $b$'s, which is all least squares needs.

How high should the degree go? Take $10$ noisy points from a wave and fit degrees $1$,
$3$ and $9$.

\begin{figure}[H]
\centering
\begin{tikzpicture}
\begin{axis}[
  name=p1, width=4.9cm, height=4.6cm, title={degree 1}, title style={font=\small},
  xmin=0, xmax=1, ymin=-2, ymax=2, xtick={0,0.5,1}, ytick={-2,-1,0,1,2},
  tick label style={font=\footnotesize}, xlabel={$x$}, ylabel={$y$},
]
\addplot[thick, sourceblue] coordinates {(0.020,0.734) (0.980,-0.718)};
\addplot[only marks, mark=*, mark size=1.6pt, densityred] coordinates {(0.05,0.11) (0.15,0.48) (0.25,0.94) (0.35,0.91) (0.45,0.59) (0.55,-0.28) (0.65,-0.95) (0.75,-1.2) (0.85,-0.62) (0.95,0.1)};
\end{axis}
\begin{axis}[
  name=p3, at={(p1.east)}, anchor=west, xshift=0.5cm,
  width=4.9cm, height=4.6cm, title={degree 3}, title style={font=\small},
  xmin=0, xmax=1, ymin=-2, ymax=2, xtick={0,0.5,1}, yticklabels={},
  tick label style={font=\footnotesize}, xlabel={$x$},
]
\addplot[thick, targetgreen, smooth] coordinates {(0.020,-0.496) (0.053,-0.048) (0.085,0.313) (0.118,0.592) (0.150,0.798) (0.183,0.935) (0.215,1.010) (0.248,1.030) (0.280,1.001) (0.313,0.929) (0.345,0.821) (0.378,0.682) (0.411,0.520) (0.443,0.341) (0.476,0.150) (0.508,-0.045) (0.541,-0.238) (0.573,-0.423) (0.606,-0.594) (0.638,-0.745) (0.671,-0.868) (0.703,-0.957) (0.736,-1.008) (0.768,-1.011) (0.801,-0.963) (0.834,-0.856) (0.866,-0.683) (0.899,-0.439) (0.931,-0.117) (0.964,0.289) (0.980,0.525)};
\addplot[only marks, mark=*, mark size=1.6pt, densityred] coordinates {(0.05,0.11) (0.15,0.48) (0.25,0.94) (0.35,0.91) (0.45,0.59) (0.55,-0.28) (0.65,-0.95) (0.75,-1.2) (0.85,-0.62) (0.95,0.1)};
\end{axis}
\begin{axis}[
  name=p9, at={(p3.east)}, anchor=west, xshift=0.5cm,
  width=4.9cm, height=4.6cm, title={degree 9}, title style={font=\small},
  xmin=0, xmax=1, ymin=-2, ymax=2, xtick={0,0.5,1}, yticklabels={},
  tick label style={font=\footnotesize}, xlabel={$x$},
]
\addplot[thick, quotientorange, smooth] coordinates {(0.022,3.019) (0.026,2.386) (0.030,1.838) (0.034,1.367) (0.038,0.965) (0.042,0.626) (0.046,0.343) (0.050,0.110) (0.060,-0.286) (0.080,-0.516) (0.100,-0.332) (0.120,0.000) (0.140,0.335) (0.160,0.604) (0.180,0.786) (0.200,0.889) (0.220,0.934) (0.240,0.943) (0.260,0.935) (0.280,0.925) (0.300,0.918) (0.320,0.916) (0.340,0.913) (0.360,0.904) (0.380,0.879) (0.400,0.832) (0.420,0.758) (0.440,0.653) (0.460,0.520) (0.480,0.362) (0.500,0.185) (0.520,-0.001) (0.540,-0.189) (0.560,-0.368) (0.580,-0.534) (0.600,-0.679) (0.620,-0.804) (0.640,-0.906) (0.660,-0.990) (0.680,-1.058) (0.700,-1.113) (0.720,-1.158) (0.740,-1.190) (0.760,-1.203) (0.780,-1.185) (0.800,-1.118) (0.820,-0.983) (0.840,-0.764) (0.860,-0.455) (0.880,-0.078) (0.900,0.303) (0.920,0.551) (0.940,0.422) (0.950,0.100) (0.954,-0.097) (0.958,-0.340) (0.962,-0.635) (0.966,-0.988) (0.970,-1.404) (0.974,-1.892) (0.978,-2.457)};
\addplot[only marks, mark=*, mark size=1.6pt, densityred] coordinates {(0.05,0.11) (0.15,0.48) (0.25,0.94) (0.35,0.91) (0.45,0.59) (0.55,-0.28) (0.65,-0.95) (0.75,-1.2) (0.85,-0.62) (0.95,0.1)};
\end{axis}
\end{tikzpicture}
\caption{The same $10$ points fitted with polynomials of degree $1$, $3$ and $9$.
Degree $1$ misses the wave (underfitting). Degree $3$ follows it. Degree $9$ has $10$
coefficients for $10$ points, so it passes through every point exactly. It wiggles to do
so, and just outside the outermost points it shoots off the chart: at $x=0.98$ it
predicts $-2.8$ (overfitting).}
\label{fig:mlr-poly}
\end{figure}

The degree-$9$ fit has zero training error, yet it predicts new data worse than any
degree from $3$ to $8$. To see this, draw $200$ fresh points from the same process and
measure mean squared error (MSE) on them, the \textbf{test error}:

\begin{figure}[H]
\centering
\begin{tikzpicture}
\begin{axis}[
  width=10.5cm, height=6cm,
  xlabel={polynomial degree}, ylabel={mean squared error},
  xmin=-0.3, xmax=9.3, ymin=0, ymax=0.7, xtick={0,...,9},
  legend style={font=\footnotesize, at={(0.98,0.98)}, anchor=north east},
  axis lines=left,
]
\addplot[thick, mark=*, sourceblue] coordinates {(0,0.5117) (1,0.3232) (2,0.3232) (3,0.0323) (4,0.0225) (5,0.0026) (6,0.0010) (7,0.0009) (8,0.0009) (9,0.0)};
\addlegendentry{training error (10 points)}
\addplot[thick, mark=square*, densityred] coordinates {(0,0.6246) (1,0.2793) (2,0.2791) (3,0.0819) (4,0.0889) (5,0.1177) (6,0.1193) (7,0.1215) (8,0.1205) (9,0.1852)};
\addlegendentry{test error (200 new points)}
\draw[gray, dashed] (axis cs:3,0) -- (axis cs:3,0.36);
\node[gray!70!black, font=\footnotesize, anchor=south west, align=left] at (axis cs:3.1,0.22) {lowest test error\\(degree $3$)};
\draw[{Stealth}-{Stealth}, gray!70!black] (axis cs:9,0.01) -- (axis cs:9,0.175);
\node[gray!70!black, font=\footnotesize, anchor=east, align=right] at (axis cs:8.85,0.27) {overfitting gap\\at degree $9$};
\end{axis}
\end{tikzpicture}
\caption{Training error never rises as the model grows; at degree $9$ it is exactly zero.
Test error falls, bottoms out at degree $3$ ($0.082$), then rises to $0.185$ at degree $9$.
The gap between the curves is overfitting: the model memorises the noise in its own
sample.}
\label{fig:mlr-traintest}
\end{figure}

\begin{connectionbox}
This U-shaped test curve is the bias--variance trade-off. Simple models are biased (they
cannot bend enough); flexible models have high variance (they bend to every quirk of the
sample). Chapter~15 develops it with train/validation/test splits, cross-validation and
regularisation (ridge and lasso), which add a penalty to $\norm{y - X\beta}^2$ to keep
coefficients small.
\end{connectionbox}

\section{Multicollinearity}
\label{sec:mlr-collinear}

Suppose we add house size measured in square feet to a model that already has size in
square metres. The two columns carry the same information: their correlation in our
$40$ houses is $0.9996$ (the square-foot figure was rounded with a little recording
noise). Here is what happens to the size coefficient:

\begin{center}
\begin{tabular}{lrr}
\toprule
model & size (m$^2$) coef & its SE\\
\midrule
size + age & 1.861 & 0.081\\
size + size\_ft$^2$ + age & 1.393 & 2.918\\
\bottomrule
\end{tabular}
\end{center}

The standard error grew about $36$ times. The fit itself barely changed ($R^2$ is
$0.938$ either way), and the combined effect $1.393 + 10.764\times 0.0432 = 1.858$ per m$^2$
is still right. The model just cannot tell how to split the credit between two
near-identical columns, so any split is roughly as good as any other.

\textbf{Multicollinearity} means some predictor is nearly a linear combination of the
others. In the projection picture, two column arrows point in almost the same direction.
The plane they span is well defined, so $\hat y$ is stable, but the coordinates of
$\hat y$ along two nearly parallel arrows are not: a tiny wobble in $y$ sends them flying
in opposite directions (Figure~\ref{fig:mlr-collinear-geom}).

\begin{figure}[H]
\centering
\begin{tikzpicture}[x=4.0cm, y=4.0cm, font=\small]
\draw[-{Stealth}, very thick, gray!70!black] (0,0) -- (1,0) node[below, font=\footnotesize] {column $x_1$};
\draw[-{Stealth}, very thick, gray!70!black] (0,0) -- (0.966,0.259) node[above, font=\footnotesize] {column $x_2$};
\draw[dashed, thick, quotientorange] (0,0) -- (0.980,0) -- (2.1,0.30);
\draw[dashed, thick, structurepurple] (0,0) -- (0.420,0) -- (2.1,0.45);
\fill[quotientorange] (0.980,0) circle (1.6pt);
\fill[structurepurple] (0.420,0) circle (1.6pt);
\draw[-{Stealth}, very thick, targetgreen] (0,0) -- (2.1,0.30);
\draw[-{Stealth}, very thick, targetgreen!60!black] (0,0) -- (2.1,0.45);
\node[anchor=west, font=\footnotesize, quotientorange] at (2.12,0.29) {$\hat y_A = 0.98\,x_1 + 1.16\,x_2$};
\node[anchor=west, font=\footnotesize, structurepurple] at (2.12,0.46) {$\hat y_B = 0.42\,x_1 + 1.74\,x_2$};
\end{tikzpicture}
\caption{Why near-duplicate columns make coefficients unstable. The grey arrows are two
predictor columns pointing almost the same way ($15^\circ$ apart). To reach a fitted
vector you walk some distance along $x_1$ (to the dot), then parallel to $x_2$ (dashed
path); those two distances are the coefficients. Two samples give green arrows
$\hat y_A$ and $\hat y_B$ that differ only slightly, yet the $x_1$ coefficient drops from
$0.98$ to $0.42$ while the $x_2$ coefficient rises from $1.16$ to $1.74$.}
\label{fig:mlr-collinear-geom}
\end{figure}

\begin{definition}[Variance inflation factor]
Regress predictor $x_j$ on all the other predictors and call the resulting $R^2$ value
$R_j^2$. The \textbf{variance inflation factor} is
\[
  \text{VIF}_j = \frac{1}{1 - R_j^2}.
\]
\end{definition}

The name is literal. One can show
\[
  \Var(\hat\beta_j) = \frac{\sigma^2}{\sum_i (x_{ij} - \bar x_j)^2}\cdot \text{VIF}_j ,
\]
so the first factor is the simple-regression variance from Chapter~12, and VIF is how
much the overlap with other predictors multiplies it. The standard error is multiplied
by $\sqrt{\text{VIF}_j}$. In the square-feet example $\text{VIF}=1274$ and
$\sqrt{1274}\approx 35.7$, matching the $36$-fold jump.

\begin{figure}[H]
\centering
\begin{tikzpicture}
\begin{axis}[
  width=10cm, height=6cm,
  xlabel={$R_j^2$ (how well other predictors explain $x_j$)},
  ylabel={SE multiplier $\sqrt{\mathrm{VIF}_j}$},
  xmin=0, xmax=1, ymin=0, ymax=8, domain=0:0.985, samples=120,
  axis lines=left,
]
\addplot[thick, structurepurple] {sqrt(1/(1-x))};
\addplot[only marks, mark=*, densityred] coordinates {(0.8,2.236) (0.9,3.162) (0.95,4.472)};
\node[anchor=east, font=\footnotesize] at (axis cs:0.78,2.4) {VIF $5$: SE $\times 2.2$};
\node[anchor=east, font=\footnotesize] at (axis cs:0.88,3.35) {VIF $10$: SE $\times 3.2$};
\node[anchor=east, font=\footnotesize] at (axis cs:0.93,4.7) {VIF $20$: SE $\times 4.5$};
\end{axis}
\end{tikzpicture}
\caption{Standard errors stay modest until a predictor is largely explained by the
others, then blow up. Common rules of thumb flag VIF above $5$ or $10$.}
\label{fig:mlr-vif}
\end{figure}

\begin{intuitionbox}
Multicollinearity hurts interpretation, not prediction. If you only want $\hat y$, a
model with two copies of size predicts fine. If you want to say ``size matters, and
by this much'', drop one of the duplicates, combine them, or collect data where they
vary separately. Perfect collinearity (VIF $=\infty$), such as the dummy trap, makes
$X^\top X$ non-invertible.
\end{intuitionbox}

\section{Adjusted \texorpdfstring{$R^2$}{R squared}}
\label{sec:mlr-adjr2}

$R^2 = 1 - \text{SSE}/\text{SST}$ never decreases when you add a column. So it cannot
tell you whether the new column helps. \textbf{Adjusted $R^2$} charges for each
coefficient:
\[
  R^2_{\text{adj}} = 1 - (1 - R^2)\,\frac{n-1}{n-p-1}.
\]
Rearranged, $1 - R^2_{\text{adj}} = \dfrac{\text{SSE}/(n-p-1)}{\text{SST}/(n-1)}$: adjusted
$R^2$ compares the residual variance to the total variance, each divided by its degrees
of freedom instead of by $n$. It is the same $n-1$ logic as the sample variance in
Chapter~2: each fitted coefficient uses up one degree of freedom. A new column lowers SSE
but also lowers $n-p-1$, so adjusted $R^2$ rises only if SSE drops by enough to pay for it.

For the $40$ houses with size, age and distance to the city centre, $R^2 = 0.9388$ and
$p = 3$:
\[
  R^2_{\text{adj}} = 1 - 0.0612\times\frac{39}{36} = 0.9337 .
\]
Drop distance and refit: $R^2$ falls slightly to $0.9377$, but with $p = 2$,
\[
  R^2_{\text{adj}} = 1 - 0.0623\times\frac{39}{37} = 0.9343,
\]
which is higher. Distance added less than the price of one coefficient.

\section{Uncertainty in the coefficients}
\label{sec:mlr-se}

Each $\hat\beta_j$ is computed from a random sample, so it has a sampling distribution
(Chapter~6). Its standard deviation is the standard error we need for confidence
intervals and tests.

\begin{theorem}[Covariance of the least squares estimator]
If $y = X\beta + \eps$ with $X$ fixed, $\E[\eps] = 0$ and $\Cov(\eps) = \sigma^2 I$
(errors uncorrelated with equal variance $\sigma^2$), then $\E[\hat\beta] = \beta$ and
\[
  \Cov(\hat\beta) = \sigma^2 (X^\top X)^{-1}.
\]
\end{theorem}

\begin{whybox}
Write $A = (X^\top X)^{-1}X^\top$, so $\hat\beta = Ay$. Substitute $y = X\beta + \eps$:
\[
  \hat\beta = AX\beta + A\eps = \beta + A\eps,
\]
since $AX = (X^\top X)^{-1}X^\top X = I$. The error term has mean zero, so $\hat\beta$
is unbiased. For a constant matrix $A$, $\Cov(A\eps) = A\,\Cov(\eps)\,A^\top$ (the
matrix version of $\Var(aX) = a^2\Var(X)$). So
\[
  \Cov(\hat\beta) = A(\sigma^2 I)A^\top = \sigma^2 (X^\top X)^{-1}X^\top X (X^\top X)^{-1}
  = \sigma^2 (X^\top X)^{-1}.
\]
\end{whybox}

The diagonal of $\Cov(\hat\beta)$ holds the variances of the individual coefficients;
the off-diagonal entries say how their errors move together. We do not know $\sigma^2$,
so we estimate it from the residuals, dividing by the residual degrees of freedom:
\[
  s^2 = \frac{\text{SSE}}{n - p - 1}, \qquad
  \SE(\hat\beta_j) = s\sqrt{\big[(X^\top X)^{-1}\big]_{jj}} .
\]

\begin{example}[Standard errors for the four houses]
Residuals were $(0.1, 0.7, -1.7, 0.9)$, so $\text{SSE} = 0.01 + 0.49 + 2.89 + 0.81 = 4.2$
and $s^2 = 4.2/(4-2) = 2.1$. The inverse was
$(X^\top X)^{-1} = \tfrac{1}{20}\begin{bmatrix}30&-10\\-10&4\end{bmatrix}$. Hence
\[
  \SE(\hat\beta_1) = \sqrt{2.1\times\tfrac{4}{20}} = \sqrt{0.42} = 0.648, \qquad
  \SE(\hat\beta_0) = \sqrt{2.1\times\tfrac{30}{20}} = \sqrt{3.15} = 1.775 .
\]
The slope SE matches Chapter~12's $s/\sqrt{S_{xx}} = \sqrt{2.1/5} = 0.648$. The
negative off-diagonal entry says that if the slope comes out too high, the intercept
tends to come out too low: the line pivots around $(\xbar,\ybar)$.
\end{example}

With normal errors, $(\hat\beta_j - \beta_j)/\SE(\hat\beta_j)$ follows a $t$
distribution with $n - p - 1$ degrees of freedom. That gives the usual recipes from
Chapters~7 and 9:
\[
  \hat\beta_j \pm t^*_{n-p-1}\,\SE(\hat\beta_j), \qquad
  t = \frac{\hat\beta_j}{\SE(\hat\beta_j)}\ \text{ for } H_0\colon \beta_j = 0 .
\]

\section{Reading a regression output table}
\label{sec:mlr-output}

Here is the full fit of price on size, age and distance (km to the centre) for the $40$
houses, in the layout printed by Python's \texttt{statsmodels}.

\begin{lstlisting}[language=Python]
import statsmodels.formula.api as smf
fit = smf.ols("price ~ size + age + dist", data=houses).fit()
print(fit.summary())
# categorical and interaction terms use the same formula language:
# smf.ols("latency ~ C(region) + payload * new_version", data=df)
\end{lstlisting}

\begin{lstlisting}[basicstyle=\ttfamily\footnotesize, numbers=none, language={}]
                            OLS Regression Results
==============================================================================
Dep. Variable:                  price   R-squared:                       0.939
Model:                            OLS   Adj. R-squared:                  0.934
Method:                 Least Squares   F-statistic:                     183.9
No. Observations:                  40   Prob (F-statistic):           6.98e-22
Df Residuals:                      36   Log-Likelihood:                -179.55
Df Model:                           3   AIC:                             367.1
                                        BIC:                             373.9
==============================================================================
                 coef    std err          t      P>|t|      [0.025      0.975]
------------------------------------------------------------------------------
Intercept     68.2558     15.856      4.305      0.000      36.099     100.413
size           1.8437      0.084     21.885      0.000       1.673       2.015
age           -1.4069      0.215     -6.547      0.000      -1.843      -0.971
dist          -0.4276      0.536     -0.797      0.431      -1.515       0.660
==============================================================================
\end{lstlisting}

Line by line, top block:
\begin{itemize}
  \item \textbf{No.\ Observations} $n = 40$. \textbf{Df Model} $p = 3$ predictors.
  \textbf{Df Residuals} $n - p - 1 = 36$: the divisor in $s^2$ and the degrees of
  freedom for every $t$ below.
  \item \textbf{R-squared} $0.939$: the model explains $93.9\%$ of the variance in
  price. \textbf{Adj.\ R-squared} $0.934$ is the penalised version from
  Section~\ref{sec:mlr-adjr2}.
  \item \textbf{F-statistic} $183.9$ and \textbf{Prob (F-statistic)} $7\times10^{-22}$:
  a test of $H_0$: \emph{all three slopes are zero}. It compares explained variance per
  predictor to residual variance per degree of freedom,
  $F = \frac{(\text{SST}-\text{SSE})/p}{\text{SSE}/(n-p-1)}$. A tiny p-value says the
  model as a whole beats predicting $\ybar$ for every house.
  \item \textbf{Log-Likelihood}, \textbf{AIC}, \textbf{BIC}: measures of fit based on
  the likelihood, with AIC and BIC adding a penalty per parameter. Lower AIC/BIC is
  better when comparing models on the same data. Chapter~14 explains likelihood.
\end{itemize}
Coefficient block, one row per column of $X$:
\begin{itemize}
  \item \textbf{coef}: $\hat\beta_j$. Holding age and distance fixed, each extra m$^2$
  adds \$1{,}844; holding size and distance fixed, each year of age subtracts \$1{,}407.
  The intercept ($68.3$) is the predicted price of a $0$ m$^2$, new house at the city
  centre: a mathematical anchor, not a real house.
  \item \textbf{std err}: $s\sqrt{[(X^\top X)^{-1}]_{jj}}$.
  \item \textbf{t}: coef divided by std err. For size, $1.8437/0.0842 = 21.9$: the
  estimate sits $22$ standard errors from zero.
  \item \textbf{P$>|$t$|$}: two-sided p-value for $H_0\colon\beta_j = 0$, from the $t_{36}$
  distribution. For distance it is $0.431$: the data are consistent with distance
  having no effect once size and age are known.
  \item \textbf{[0.025\ \ 0.975]}: the $95\%$ confidence interval
  $\hat\beta_j \pm 2.028\,\SE$, since $t^*_{36} = 2.028$. For size:
  $1.8437 \pm 2.028\times 0.0842 = [1.673,\ 2.015]$. The distance interval
  $[-1.515,\ 0.660]$ contains $0$, matching its large p-value.
\end{itemize}
The residual standard error, $s = \sqrt{\text{SSE}/36} = 22.7$, is not printed in the
summary. In \texttt{statsmodels} it is the square root of \texttt{fit.scale}. It says
typical prediction misses are around \$23k.

\begin{figure}[H]
\centering
\begin{tikzpicture}
\begin{axis}[
  width=10.5cm, height=5.8cm,
  xlabel={fitted price $\hat y$ (\$k)}, ylabel={residual $e$ (\$k)},
  xmin=60, xmax=420, ymin=-65, ymax=65, axis lines=left,
]
\addplot[only marks, mark=*, mark size=1.7pt, sourceblue] coordinates {(308.8,-0.8) (324.8,18.2) (324.6,-11.6) (149.1,11.9) (180.4,-38.4) (331.6,-17.6) (149.0,4.0) (337.9,-50.9) (334.8,23.2) (213.5,-29.5) (211.6,18.4) (177.8,-9.8) (219.7,12.3) (245.6,4.4) (270.1,-40.1) (292.7,26.3) (365.5,44.5) (343.2,-3.2) (248.2,8.8) (382.4,-2.4) (158.4,-12.4) (144.0,39.0) (267.5,-4.5) (154.9,-3.9) (123.7,-12.7) (278.7,-20.7) (249.2,-30.2) (399.0,22.0) (248.1,12.9) (278.3,19.7) (234.1,10.9) (193.6,-15.6) (89.0,9.0) (147.6,0.4) (340.5,-10.5) (135.8,9.2) (172.9,13.1) (84.5,-16.5) (338.3,-16.3) (181.4,39.6)};
\draw[densityred, thick] (axis cs:60,0) -- (axis cs:420,0);
\draw[gray, dashed] (axis cs:60,45.4) -- (axis cs:420,45.4);
\draw[gray, dashed] (axis cs:60,-45.4) -- (axis cs:420,-45.4);
\node[gray!70!black, font=\footnotesize, anchor=south east] at (axis cs:418,45.4) {$+2s = +45$};
\node[gray!70!black, font=\footnotesize, anchor=north east] at (axis cs:418,-45.4) {$-2s = -45$};
\end{axis}
\end{tikzpicture}
\caption{Before trusting the table, check the residuals against fitted values. Here they
form a patternless band around zero with roughly constant spread, and nearly all lie
within the dashed lines at $\pm 2s = \pm 45$ (one sits just on the upper line). A funnel or a curve here would mean the standard errors in the
table are not to be trusted.}
\label{fig:mlr-resid}
\end{figure}

\begin{trapbox}
A small p-value says a coefficient is distinguishable from zero, not that it is large or
causal. A large p-value (distance, $0.43$) does not prove the effect is zero: the
interval runs from $-1.5$ to $+0.7$ thousand dollars per km, which includes effects
someone might care about. And the p-values are computed one at a time; scanning twenty
of them for stars runs into the multiple-testing problem of Chapter~11.
\end{trapbox}

\section{Logistic regression: predicting yes or no}
\label{sec:logit}

\subsection{Why a straight line fails}

Twenty visitors to a product page: minutes spent on the page and whether they bought
($1$) or not ($0$).

\begin{center}
\begin{tabular}{l*{10}{r}}
\toprule
minutes & 1 & 2 & 2 & 3 & 4 & 4 & 5 & 5 & 6 & 6\\
bought  & 0 & 0 & 0 & 0 & 0 & 1 & 0 & 0 & 1 & 0\\
\midrule
minutes & 7 & 7 & 8 & 8 & 9 & 10 & 11 & 12 & 13 & 14\\
bought  & 1 & 0 & 1 & 1 & 0 & 1 & 1 & 1 & 1 & 1\\
\bottomrule
\end{tabular}
\end{center}

We want $\P(\text{buy}\given\text{minutes})$. Fitting an ordinary line to the $0$/$1$
column gives $\hat y = -0.120 + 0.0906\cdot\text{minutes}$. At $0$ minutes it predicts a
probability of $-0.12$; at $16$ minutes, $1.33$. Neither is a probability. The errors are
not normal either (each $y$ is $0$ or $1$) and their variance $p(1-p)$ changes with $x$,
so the standard errors from Section~\ref{sec:mlr-se} are wrong too.

\begin{figure}[H]
\centering
\begin{tikzpicture}
\begin{axis}[
  width=11cm, height=6.3cm,
  xlabel={minutes on page}, ylabel={$\P(\text{buy})$},
  xmin=0, xmax=16, ymin=-0.25, ymax=1.4, ytick={0,0.5,1},
  legend style={font=\footnotesize, at={(0.98,0.2)}, anchor=south east},
  axis lines=left,
]
\fill[trapred!12] (axis cs:0,1) rectangle (axis cs:16,1.4);
\fill[trapred!12] (axis cs:0,-0.25) rectangle (axis cs:16,0);
\addplot[thick, dashed, gray!70!black, domain=0:16] {-0.1203+0.09056*x};
\addlegendentry{straight line}
\addplot[thick, sourceblue, domain=0:16, samples=100] {1/(1+exp(-(-4.0455+0.61032*x)))};
\addlegendentry{logistic curve}
\addplot[only marks, mark=*, mark size=1.8pt, densityred] coordinates
  {(1,0) (2,0.03) (2,-0.03) (3,0) (4,0) (5,0.03) (5,-0.03) (6,0) (7,0) (9,0)
   (4,1) (6,1) (7,1) (8,1.03) (8,0.97) (10,1) (11,1) (12,1) (13,1) (14,1)};
\addlegendentry{visitors}
\node[trapred, font=\footnotesize, anchor=north west] at (axis cs:0.2,1.38) {impossible: $p>1$};
\node[trapred, font=\footnotesize, anchor=south east] at (axis cs:15.8,-0.24) {impossible: $p<0$};
\end{axis}
\end{tikzpicture}
\caption{The $20$ visitors: dots at $1$ bought, dots at $0$ did not (two visitors with the same
minutes are nudged slightly apart). The straight line leaves $[0,1]$ into the shaded
impossible regions below $1.3$ minutes and above $12.4$ minutes. The logistic curve bends
to stay inside, approaching $0$ and $1$ without crossing.}
\label{fig:logit-vs-linear}
\end{figure}

\subsection{Odds, log-odds and the sigmoid}

The fix is to model a quantity that can take any real value, then squash it back into
$[0,1]$. Two steps get there.

The \textbf{odds} of an event with probability $p$ are $p/(1-p)$: how many times more
likely it is to happen than not. Probability $0.8$ is odds of $4$ (``4 to 1''); probability
$0.5$ is odds $1$. Odds range over $(0,\infty)$.

The \textbf{log-odds}, or \textbf{logit}, is $\log\frac{p}{1-p}$ (natural log). It
ranges over all real numbers: $p = 0.5$ gives $0$, $p = 0.8$ gives $\log 4 = 1.386$,
$p = 0.2$ gives $-1.386$. A line can live on that scale without hitting any wall.

\begin{definition}[Logistic regression]
For a $0$/$1$ response with predictors $x_1,\dots,x_p$, logistic regression models
\[
  \log\frac{p}{1-p} = \beta_0 + \beta_1x_1 + \dots + \beta_p x_p = x^\top\beta,
\]
equivalently
\[
  p = \sigma(x^\top\beta), \qquad \sigma(z) = \frac{1}{1 + e^{-z}} .
\]
The function $\sigma$ is the \textbf{sigmoid} (or logistic function): it maps any real
number $z$ to a probability.
\end{definition}

\begin{figure}[H]
\centering
\begin{tikzpicture}
\begin{axis}[
  width=11cm, height=6cm,
  xlabel={$z = x^\top\beta$ (log-odds)}, ylabel={$p = \sigma(z)$},
  xmin=-6.5, xmax=6.5, ymin=-0.05, ymax=1.12, ytick={0,0.2,0.5,0.8,1},
  xtick={-6,-4,-1.386,0,1.386,4,6}, xticklabels={$-6$,$-4$,$-1.39$,$0$,$1.39$,$4$,$6$},
  tick label style={font=\footnotesize}, axis lines=left,
]
\addplot[very thick, sourceblue, domain=-6.5:6.5, samples=120] {1/(1+exp(-x))};
\draw[gray, dashed] (axis cs:-6.5,1) -- (axis cs:6.5,1);
\draw[gray, dotted] (axis cs:1.386,0) -- (axis cs:1.386,0.8) -- (axis cs:-6.5,0.8);
\draw[gray, dotted] (axis cs:-1.386,0) -- (axis cs:-1.386,0.2) -- (axis cs:-6.5,0.2);
\fill[densityred] (axis cs:0,0.5) circle (2pt);
\node[densityred, font=\footnotesize, anchor=east] at (axis cs:-0.3,0.58) {$z=0 \Rightarrow p=0.5$};
\draw[densityred, thick, dashed] (axis cs:-2,0) -- (axis cs:2,1);
\node[densityred, font=\footnotesize, anchor=east] at (axis cs:1.2,0.9) {tangent: slope $1/4$};
\node[font=\footnotesize, anchor=north] at (axis cs:5.2,0.8) {$p\to 1$};
\end{axis}
\end{tikzpicture}
\caption{The sigmoid turns log-odds into probability. It is symmetric about $(0, 0.5)$,
steepest there (dashed tangent, slope $1/4$), and flattens toward $0$ and $1$ without
reaching them. The dotted guides read off two values: log-odds $+1.39$ is probability
$0.8$ and log-odds $-1.39$ is probability $0.2$.}
\label{fig:logit-sigmoid}
\end{figure}

\subsection{Interpreting coefficients as odds ratios}

Fitting the page data (method in the next subsection) gives
\[
  \log\frac{\hat p}{1-\hat p} = -4.045 + 0.610\cdot\text{minutes}.
\]
On the log-odds scale this is a line: each extra minute adds $0.610$ to the log-odds.
Exponentiate and it becomes multiplication: each extra minute multiplies the odds by
\[
  e^{0.610} = 1.84 .
\]
That number is the \textbf{odds ratio}: the factor by which the odds change for a one-unit
increase in the predictor, holding other predictors fixed.

\begin{center}
\begin{tabular}{rrrr}
\toprule
minutes & log-odds $-4.045 + 0.610x$ & odds & $\hat p$\\
\midrule
2 & $-2.825$ & 0.059 & 0.056\\
5 & $-0.994$ & 0.370 & 0.270\\
8 & $0.837$ & 2.31 & 0.698\\
12 & $3.278$ & 26.5 & 0.964\\
\bottomrule
\end{tabular}
\end{center}

From $5$ to $8$ minutes the odds go from $0.370$ to $2.31$, a factor of
$2.31/0.370 = 6.24 = 1.84^3$. The probability rises by $0.43$. From $8$ to $11$ minutes
the odds again grow $6.24$-fold, but the probability can only rise from $0.70$ to about
$0.94$. Constant effect on the odds; diminishing effect on the probability
(Figure~\ref{fig:logit-oddsratio}).

\begin{figure}[H]
\centering
\begin{tikzpicture}
\begin{axis}[
  width=11cm, height=6.3cm,
  xlabel={minutes on page}, ylabel={$\hat p$ = predicted $\P(\text{buy})$},
  xmin=0, xmax=16, ymin=0, ymax=1.08, ytick={0,0.27,0.70,0.94,1},
  yticklabels={$0$,$0.27$,$0.70$,$0.94$,$1$}, xtick={0,2,5,8,11,14,16},
  tick label style={font=\footnotesize}, axis lines=left,
]
\addplot[thick, sourceblue, domain=0:16, samples=120] {1/(1+exp(-(-4.0455+0.61032*x)))};
\draw[gray, dotted] (axis cs:0,0.270) -- (axis cs:5,0.270);
\draw[gray, dotted] (axis cs:0,0.698) -- (axis cs:8,0.698);
\draw[gray, dotted] (axis cs:0,0.935) -- (axis cs:11,0.935);
\addplot[only marks, mark=*, mark size=2pt, densityred] coordinates {(5,0.270) (8,0.698) (11,0.935)};
\draw[-{Stealth}, thick, densityred] (axis cs:8.4,0.29) -- (axis cs:8.4,0.68);
\node[densityred, font=\footnotesize, anchor=west, align=left] at (axis cs:8.5,0.42) {$5\to8$ min:\\odds $\times 6.24$\\$\hat p$ rises $0.43$};
\draw[-{Stealth}, thick, quotientorange] (axis cs:11.4,0.71) -- (axis cs:11.4,0.92);
\node[quotientorange, font=\footnotesize, anchor=north west, align=left] at (axis cs:11.5,0.86) {$8\to11$ min:\\odds $\times 6.24$\\$\hat p$ rises only $0.24$};
\end{axis}
\end{tikzpicture}
\caption{The fitted curve for the page data. Each three-minute step multiplies the odds of
buying by the same $1.84^3 = 6.24$, but the probability gained shrinks as $\hat p$ nears
$1$: $+0.43$ from $5$ to $8$ minutes, only $+0.24$ from $8$ to $11$.}
\label{fig:logit-oddsratio}
\end{figure}

\begin{trapbox}
An odds ratio of $1.84$ does not mean ``$84\%$ more likely to buy''. It means the odds,
not the probability, grow by $84\%$. When $p$ is small the two are close (odds
$\approx p$), which is why the confusion survives; when $p$ is near $0.5$ or above they
diverge sharply.
\end{trapbox}

\begin{recoverybox}
Pause and think: a coefficient of $\beta = 0$ gives an odds ratio of $e^0 = 1$. What does
a negative coefficient do to the odds? (It gives an odds ratio below $1$: $\beta = -0.69$
halves the odds.)
\end{recoverybox}

\subsection{Fitting by maximum likelihood}

There are no residuals to square in the usual sense, so we fit differently. For each
candidate $\beta$, the model assigns visitor $i$ a probability $p_i = \sigma(x_i^\top\beta)$
of buying. The probability the model assigns to what \emph{actually happened} is $p_i$
if they bought and $1 - p_i$ if not, which can be written in one expression as
$p_i^{y_i}(1-p_i)^{1-y_i}$. Assuming visitors are independent, multiply over all of them:
\[
  L(\beta) = \prod_{i=1}^n p_i^{\,y_i}(1-p_i)^{1-y_i}.
\]
$L$ is the \textbf{likelihood}: how probable the observed data are under a given
$\beta$. \textbf{Maximum likelihood} picks the $\beta$ that makes the data most probable.
Products of many small numbers are awkward, so take logs and flip the sign to get a loss
to minimise:
\[
  \ell(\beta) = -\sum_{i=1}^n \Big[\,y_i\log p_i + (1-y_i)\log(1-p_i)\Big].
\]
Divided by $n$, this is the \textbf{log loss}, also called binary cross-entropy. Each
visitor contributes $-\log(\text{probability given to the true outcome})$.

\begin{figure}[H]
\centering
\begin{tikzpicture}
\begin{axis}[
  width=11.5cm, height=6cm,
  xlabel={predicted probability of buying $p$}, ylabel={loss for one visitor},
  xmin=0, xmax=1, ymin=0, ymax=4.6, domain=0.01:0.99, samples=150,
  axis lines=left,
]
\addplot[thick, sourceblue] {-ln(x)};
\addplot[thick, densityred, dashed] {-ln(1-x)};
\node[sourceblue, font=\footnotesize, anchor=west, align=left] at (axis cs:0.1,3.9) {visitor bought ($y=1$):\\loss $-\log p$};
\node[densityred, font=\footnotesize, anchor=east, align=right] at (axis cs:0.94,3.9) {visitor did not buy ($y=0$):\\loss $-\log(1-p)$};
\addplot[only marks, mark=*, mark size=2pt, sourceblue] coordinates {(0.9,0.105) (0.1,2.303)};
\node[sourceblue, font=\footnotesize, anchor=south west] at (axis cs:0.9,0.15) {$0.11$};
\node[sourceblue, font=\footnotesize, anchor=west, align=left] at (axis cs:0.15,2.5) {buyer given $p=0.1$:\\loss $2.30$};
\end{axis}
\end{tikzpicture}
\caption{Log loss punishes confident mistakes hardest. For a buyer (solid curve), predicting
$p = 0.9$ costs only $0.11$, but predicting $p = 0.1$ costs $2.30$ (marked dots). Each
curve shoots to infinity as the prediction becomes certain and wrong.}
\label{fig:logit-logloss}
\end{figure}

There is no closed-form solution like the normal equations. But $\ell$ is convex (a
single bowl with no false bottoms), so iterative methods find the minimum reliably.
Software uses Newton's method, which re-solves a weighted least squares problem at each
step; machine learning libraries use gradient descent. Both use the same gradient,
\[
  \nabla\ell(\beta) = X^\top(p - y),
\]
which has the same shape as the least squares gradient $X^\top(X\beta - y)$, with
predicted probabilities in place of predicted values. At the optimum the residuals
$y - \hat p$ are again orthogonal to every column of $X$.

For the page data the minimum log loss is $0.415$ per visitor. Standard errors come from
the curvature of $\ell$ at its minimum: the analogue of $\sigma^2(X^\top X)^{-1}$ is
$(X^\top W X)^{-1}$, where $W$ is diagonal with entries $\hat p_i(1-\hat p_i)$. Here
$\SE(\hat\beta_1) = 0.268$, so $z = 0.610/0.268 = 2.27$ and $p = 0.023$. The $95\%$
interval for $\beta_1$ is $0.610 \pm 1.96\times 0.268$; exponentiating its endpoints gives
an odds-ratio interval of $[1.09,\ 3.12]$.

\begin{lstlisting}[language=Python]
import numpy as np, statsmodels.formula.api as smf
fit = smf.logit("bought ~ minutes", data=visits).fit()
print(fit.params)                 # Intercept -4.045, minutes 0.610
print(np.exp(fit.params))         # odds ratio per minute: 1.84
print(np.exp(fit.conf_int()))     # odds-ratio CI: [1.09, 3.12]
# sklearn.linear_model.LogisticRegression adds a ridge penalty by default
\end{lstlisting}

\begin{trapbox}
If some combination of predictors separates the $0$s from the $1$s perfectly, the
likelihood keeps increasing as the coefficients grow without bound, and the fit
``converges'' to huge numbers with huge standard errors. This is called
\textbf{complete separation}. The fix is more data, fewer predictors, or a penalty on
coefficient size.
\end{trapbox}

\subsection{From probabilities to decisions}

A classifier needs a yes/no answer. The usual rule is: predict $1$ when $\hat p \ge 0.5$.
Since $\sigma(z) \ge 0.5$ exactly when $z \ge 0$, this is the same as
$x^\top\hat\beta \ge 0$. For the page data, $-4.045 + 0.610x \ge 0$ when
$x \ge 6.63$ minutes.

With two predictors, $\hat\beta_0 + \hat\beta_1 x_1 + \hat\beta_2 x_2 = 0$ is a straight
line in the $(x_1, x_2)$ plane: the \textbf{decision boundary}. With $p$ predictors it is a
flat hyperplane. Figure~\ref{fig:logit-boundary} shows $30$ users described by two
engagement scores, with the fitted model
$\log\frac{\hat p}{1-\hat p} = -12.55 + 1.94x_1 + 1.98x_2$.

\begin{figure}[H]
\centering
\begin{tikzpicture}
\begin{axis}[
  width=9.5cm, height=8cm,
  xlabel={score $x_1$}, ylabel={score $x_2$},
  xmin=0, xmax=6, ymin=0, ymax=6, axis lines=left,
  legend style={font=\footnotesize, at={(0.5,-0.17)}, anchor=north, legend columns=2,
    /tikz/every even column/.append style={column sep=0.5cm}},
]
\addplot[thick, black, domain=0:6] {(12.55179-1.94243*x)/1.98008};
\addlegendentry{$\hat p = 0.5$ boundary}
\addplot[thin, dashed, gray, domain=0:6] {(12.55179+1.38629-1.94243*x)/1.98008};
\addlegendentry{$\hat p = 0.2$ and $0.8$}
\addplot[thin, dashed, gray, domain=0:6, forget plot] {(12.55179-1.38629-1.94243*x)/1.98008};
\addplot[only marks, mark=o, mark size=2.3pt, thick, sourceblue] coordinates {(2.6,3.7) (0.4,3.9) (2.2,1.4) (1.4,1.1) (2.0,3.0) (2.8,2.8) (0.5,0.6) (3.8,3.2) (4.4,3.4) (1.2,3.5) (2.8,2.4) (1.4,2.2) (2.1,3.1) (2.3,3.8) (1.4,2.7)};
\addlegendentry{did not convert}
\addplot[only marks, mark=triangle*, mark size=2.8pt, densityred] coordinates {(4.3,2.9) (2.6,4.8) (4.3,3.6) (4.2,3.1) (4.0,3.2) (4.0,4.8) (3.5,4.6) (2.8,3.5) (2.4,3.7) (3.9,3.6) (3.5,5.1) (5.2,2.6) (5.0,2.1) (3.1,3.6) (3.9,2.8)};
\addlegendentry{converted}
\addplot[only marks, mark=o, mark size=5pt, thick, quotientorange] coordinates {(3.8,3.2) (4.4,3.4) (2.8,3.5) (2.4,3.7)};
\addlegendentry{misclassified (wrong side)}
\node[font=\footnotesize, anchor=south west, sourceblue] at (axis cs:1.6,0.2) {predict $0$};
\node[font=\footnotesize, anchor=north east, densityred] at (axis cs:5.8,5.8) {predict $1$};
\end{axis}
\end{tikzpicture}
\caption{The decision boundary of logistic regression is a straight line where the
log-odds equal zero: above-right the model predicts ``convert'', below-left ``not''. The
dashed lines are where $\hat p = 0.8$ (upper) and $\hat p = 0.2$ (lower). The four ringed
users fall on the wrong side, three of them inside that uncertain band, so training
accuracy is $26/30$.}
\label{fig:logit-boundary}
\end{figure}

\begin{connectionbox}
A logistic regression is a neural network with no hidden layer: a weighted sum of inputs
followed by a sigmoid, trained by minimising cross-entropy with gradient descent. Replace
the sigmoid by a softmax over $K$ classes and you get multinomial logistic regression,
the output layer of nearly every classifier. Curved boundaries come from feeding in
nonlinear features, exactly as polynomial columns bent the regression line.
\end{connectionbox}

\section{Generalised linear models: the umbrella}
\label{sec:glm}

Linear and logistic regression share a three-part structure:
\begin{enumerate}
  \item a \textbf{linear predictor} $\eta = x^\top\beta$, the weighted sum of inputs;
  \item a \textbf{link function} $g$ that connects the mean of $y$ to $\eta$ through
  $g(\E[y\given x]) = \eta$;
  \item a \textbf{distribution} for $y$ around its mean.
\end{enumerate}
Linear regression uses the identity link and a normal distribution. Logistic regression
uses the logit link and a Bernoulli distribution (a single $0$/$1$ trial). Counts, such as
errors per hour or clicks per page view, use a log link and a Poisson distribution, so
$\log\E[y] = x^\top\beta$ and each coefficient multiplies the expected count by
$e^{\beta_j}$. Any model of this form is a \textbf{generalised linear model} (GLM), and
all of them are fitted by maximum likelihood with the same algorithm.

\begin{figure}[H]
\centering
\begin{tikzpicture}[font=\small,
  top/.style={draw=structurepurple, very thick, rounded corners, align=center, fill=structurepurple!8, minimum width=6.5cm, minimum height=1cm},
  leaf/.style={draw, thick, rounded corners, align=center, minimum width=3.4cm, minimum height=1.9cm, text width=3.2cm}]
\node[top] (G) at (0,0) {Generalised linear model\\$g(\E[y\given x]) = x^\top\beta$};
\node[leaf, draw=sourceblue, fill=sourceblue!6] (L) at (-4,-2.6) {\textbf{Linear}\\$y$ continuous\\link: identity\\Normal errors};
\node[leaf, draw=densityred, fill=densityred!6] (B) at (0,-2.6) {\textbf{Logistic}\\$y\in\{0,1\}$\\link: logit\\Bernoulli};
\node[leaf, draw=targetgreen, fill=targetgreen!6] (P) at (4,-2.6) {\textbf{Poisson}\\$y$ a count\\link: log\\Poisson};
\draw[-{Stealth}, thick] (G) -- (L);
\draw[-{Stealth}, thick] (G) -- (B);
\draw[-{Stealth}, thick] (G) -- (P);
\node[anchor=north, font=\footnotesize, align=center] at (L.south) {coef: additive\\change in $y$};
\node[anchor=north, font=\footnotesize, align=center] at (B.south) {$e^{\text{coef}}$: odds\\ratio};
\node[anchor=north, font=\footnotesize, align=center] at (P.south) {$e^{\text{coef}}$: rate\\ratio};
\end{tikzpicture}
\caption{Generalised linear models share one linear predictor and differ in the link
function and the distribution of $y$. The link determines how to read a coefficient:
additively, as an odds ratio, or as a rate ratio.}
\label{fig:glm-umbrella}
\end{figure}

\begin{historybox}
In 1944 the American statistician Joseph Berkson, working on dose--response data (what
fraction of animals respond at each dose of a drug), proposed the log-odds transform and
named it the ``logit''. The competing approach of the day, the ``probit'', used the normal
distribution's cumulative curve and needed tables of the normal to fit; the logit had a
simple closed form. For decades, regression for continuous data, for proportions and for
counts were taught as separate subjects with separate fitting tricks. In 1972 John Nelder
and Robert Wedderburn showed they were one model with different link functions and
distributions, all fittable by a single iterative weighted least squares algorithm. That
paper created the term ``generalised linear model'' and the software (GLIM, and later R's
\texttt{glm}) that made logistic and Poisson regression routine.
\end{historybox}

\section{Summary}
\label{sec:mlr-summary}

\begin{itemize}
  \item \textbf{Matrix form}: $y = X\beta + \eps$, with design matrix $X$ (rows =
  observations, columns = intercept plus predictors).
  \item \textbf{Normal equations}: $X^\top X\hat\beta = X^\top y$, so
  $\hat\beta = (X^\top X)^{-1}X^\top y$ when the columns are independent.
  \item \textbf{Geometry}: $\hat y = Hy$ is the orthogonal projection of $y$ onto the column
  space of $X$; the residual satisfies $X^\top e = 0$. Leverages $h_{ii}$ sum to $p+1$.
  \item \textbf{Interpretation}: $\hat\beta_j$ is the change in $\hat y$ per unit of
  $x_j$ with the other predictors fixed. Adding or dropping a correlated variable can change
  it, even flip its sign (omitted variable bias).
  \item \textbf{Features}: dummies ($k-1$ for $k$ categories), interactions (products,
  non-parallel lines), polynomials (curves). More flexibility lowers training error and
  eventually raises test error.
  \item \textbf{Multicollinearity}: $\text{VIF}_j = 1/(1-R_j^2)$ multiplies
  $\Var(\hat\beta_j)$; it hurts interpretation, not prediction.
  \item \textbf{Adjusted $R^2$} $= 1-(1-R^2)\frac{n-1}{n-p-1}$ can fall when a useless
  variable is added.
  \item \textbf{Uncertainty}: $\Cov(\hat\beta) = \sigma^2(X^\top X)^{-1}$, estimated with
  $s^2 = \text{SSE}/(n-p-1)$; $t$-tests and CIs use $n-p-1$ degrees of freedom.
  \item \textbf{Logistic regression}: $\log\frac{p}{1-p} = x^\top\beta$,
  $p = \sigma(x^\top\beta)$; $e^{\beta_j}$ is an odds ratio; fit by maximum likelihood
  (minimum log loss); decision boundary $x^\top\hat\beta = 0$ is linear.
  \item \textbf{GLMs}: linear predictor + link + distribution; linear, logistic and
  Poisson regression are special cases.
\end{itemize}

Chapter~14 takes the idea that fitted logistic regression, maximum likelihood, and
develops it as a general method for estimation, together with its Bayesian counterpart.

\section{Exercises}
\label{sec:mlr-exercises}

\begin{enumerate}
  \item For $x = (0,1,2,3,4)$ and $y = (1,3,2,5,4)$, write the design matrix, compute
  $X^\top X$ and $X^\top y$, and solve the normal equations.
  \item A model gives $\widehat{\text{price}} = 50 + 2.0\cdot\text{size} - 1.5\cdot\text{age}$
  (price in \$k, size in m$^2$, age in years). Predict the price of a $120$ m$^2$,
  $10$-year-old house, and state in words what $-1.5$ means.
  \item A regression of salary on years of experience gives a slope of $3$ (\$k per year).
  Adding education changes it to $2$. Using the omitted variable bias formula, what does
  this tell you about the relationship between experience and education, given that
  education has a positive coefficient?
  \item A dataset has a ``day of week'' column with $7$ values. How many dummy columns
  should a model with an intercept use, and what goes wrong if you use $7$?
  \item In $\hat y = 10 + 2x + 5d + 3xd$ with $d\in\{0,1\}$, give the intercept and slope
  of the line for each group. Are the lines parallel?
  \item With $n = 30$, $p = 5$ and $R^2 = 0.40$, compute adjusted $R^2$.
  \item A predictor has $R_j^2 = 0.8$ when regressed on the other predictors. By what
  factor is its standard error inflated?
  \item Show that the hat matrix satisfies $HH = H$, and explain geometrically why this
  must be true.
  \item In a logistic regression for churn, the coefficient on ``received discount'' is
  $0.405$. Compute the odds ratio. If a customer's churn probability is $0.20$ without the
  discount, what is it with the discount?
  \item A logistic model predicts $\hat p = 0.8$ for one positive example and
  $\hat p = 0.3$ for one negative example. Compute the average log loss over the two.
\end{enumerate}

\subsection*{Answers}

\begin{enumerate}
  \item $X$ has rows $(1, x_i)$. $X^\top X = \begin{bmatrix}5&10\\10&30\end{bmatrix}$,
  $X^\top y = (15,\ 38)$. Determinant $50$, so
  $\hat\beta = \tfrac{1}{50}(30\cdot 15 - 10\cdot 38,\ -10\cdot 15 + 5\cdot 38) = (1.4,\ 0.8)$.
  \item $50 + 240 - 15 = 275$, so \$275k. Among houses of the same size, each extra year
  of age lowers the predicted price by \$1{,}500.
  \item Simple slope $=$ direct slope $+$ (education coef)$\times\delta$, so
  $3 = 2 + (\text{education coef})\times\delta$. The product is $1 > 0$; with a positive
  education coefficient, $\delta > 0$: more experienced people in this data also have more
  education, and part of the simple slope was education's effect.
  \item Six, with one day as baseline. With seven, the dummies sum to the intercept column,
  $X^\top X$ is singular and the coefficients are not identifiable.
  \item $d = 0$: $\hat y = 10 + 2x$. $d = 1$: $\hat y = 15 + 5x$. Not parallel; the
  interaction coefficient $3$ is the difference in slopes.
  \item $1 - 0.60\times 29/24 = 1 - 0.725 = 0.275$.
  \item $\text{VIF} = 1/0.2 = 5$; the SE is multiplied by $\sqrt5 \approx 2.24$.
  \item $HH = X(X^\top X)^{-1}X^\top X(X^\top X)^{-1}X^\top = X(X^\top X)^{-1}X^\top = H$.
  Projecting a vector that already lies in the column space leaves it where it is.
  \item Odds ratio $e^{0.405} \approx 1.50$. Odds $0.2/0.8 = 0.25$ become $0.375$, so
  $p = 0.375/1.375 \approx 0.273$.
  \item Losses $-\log 0.8 = 0.223$ and $-\log(1-0.3) = -\log 0.7 = 0.357$; average
  $0.290$.
\end{enumerate}
